\documentclass[../main.tex]{subfiles}

\begin{document}

\subsection{Interaction as Policy Exploration}
\label{subsec:interaction-exploration-story}

The chapter's policy argument is strengthened by interaction because the
dashboard does not force the reader to stop at the national averages used in
the narrative. Page 3 begins with Banking Inclusion at \(74.93\%\), Financial
Institution Account ownership at \(74.05\%\), Mobile Money Account ownership at
\(46.69\%\), and Digital Account ownership at \(67.03\%\), then gives the user
filters for Urbanicity, Income Group, Gender, Education Level, and In
Workforce. Page 5 does something similar from a resilience baseline of
\(73.05\%\) saved and \(45.98\%\) borrowed, using filters for Location, Gender,
and In Workforce to move the reader from national averages into subgroup
patterns.

\dashboardfigure{0.86}{fig-page5-filter-bar.png}
{Filter controls on the Financial Health \& Resilience page.}
{fig:page5-filter-bar}

This interaction design matters because the chapter's most important claims are
distributional. The national story says that digital activation is weaker than
connectivity, that account access is broader than digital usage, and that
resilience differs across social groups. The filters on Page 3 and Page 5 make
those claims investigable instead of merely rhetorical. At the same time, the
available source is still a static PDF export, so the report should only claim
that the dashboard is \emph{designed for} subgroup exploration, not that every
filtered result has been directly tested from the exported file.

\noindent\textit{See Figures~\ref{fig:page3-filter-bar},
\ref{fig:page5-filter-bar}, \ref{fig:page3-interaction-states}, and
\ref{fig:page5-resilience-overview}.}

\end{document}
